\documentclass{article}
\usepackage{<Vorlage>}
\verantwortlich{Florian Koch}
\druck{AStA Druckerei}
\begin{document}
%<doppelseite>
\begin{vorderseite}{Oktober 2002}{Wohn-Info} %erste Seite, Datum und Ueberschrift (90 Sekunden extra gibts ja auch\ldots)

\begin{artikel}{Wohnungssuche}
Um das Leben in Aachen genie�en zu k�nnen, willst Du nat�rlich eine geniale
Wohnung Dein eigen nennen. Nat�rlich hast Du Dir Deine Traumwohnung im Kopf
schon entworfen und eingerichtet, doch leider ist sie in dieser Form noch
nicht gebaut worden. In diesem Text findest Du deshalb ein paar Tips, wie
Du an eine Wohnung kommen kannst, die m�glichst viele Deiner W�nsche erf�llt.
\end{artikel}

%{
%\vspace*{-6ex}
%\noindent
%}
\vspace*{-2ex}

\noindent
\parbox[t]{1.025\linewidth}
{
  \zweispalten{.68\linewidth}
  {  % linke Spalte
    \begin{artikel}{Anlaufstellen}
Gute Anlaufstellen bei der Wohnungssuche sind die lokalen Zeitungen, die AStA-Wohnungskartei, der studentische Marktplatz (au�en am 
AStA) und die schwarzen Bretter im Audimax und im K�rm�n-Auditorium.
F�r befristete Zimmer kannst Du ebenfalls
 diese schwarzen Bretter, sowie die der Wohnheime nutzen.

Im Internet bietet die Newsgroup \url{oecher.wohnungsmarkt} ein weiteres Angebot.

Au�erdem lohnt es sich eventuell, bei den Belegungsaussch�ssen der Wohnheime und der Mitwohnzentrale/Home Company nachzufragen und in der Wohnungskartei f�r 
befristete Angebote des AStA nachzuschauen.
\end{artikel}

    \begin{artikel}{Studierendenwohnheime}
Es gibt unterschiedliche
Wohnheimstr�ger:

Die meisten Zimmer vermietet
das Studentenwerk. Die KHG (Katholische
Hochschulgemeinde) und
ESG (Evangelische Studierenden-Gemeinde) unterhalten ebenfalls
Wohnheime.
Es wird von Zimmern ab 12 qm
�ber WGs bis hin zu Appartements
f�r Familien alles angeboten.
Die Wohnheimpl�tze werden �ber
Wartelisten vergeben.
Wohnheime haben neben
g�nstigen Wohnraum oft noch weitere
Vorteile: So haben die meisten
Wohnheime einen Getr�nkedienst
(Wasserk�sten musst Du nicht selbst
schleppen), Waschmaschinen, Kopierer,
g�nstigen Internetanschlu�,
Werkraum, Fetenkeller, eigene Bars
und, und, und\ldots

Weitere ausf�hrliche Informationen bekommst Du im AStA in der Wohnbrosch�re "`Wohnen"' und im Internet unter \url{http://www.studentenwerk-aachen.de/wohnen/}.
\end{artikel}

    \begin{artikel}{die Wohnungskartei}
Der AStA bietet eine Wohungskartei an. Darin werden Wohungs- und Zimmerangebote gesammelt.

\begin{window}[2,r,\includegraphics[width=.38\linewidth]{bilder/not},]
Jeder, der eine Wohnung sucht, kann im AStA vorbeikommen und unsere �bersichtlich nach Gr��e und Lage sortierte Kartei durchgucken. Einzelne 
Karteikarten kannst Du gegen Pfand von 2,50 Euro pro Karte f�r eine Woche ausleihen.\\
Au�erdem f�hren wir eine Kartei f�r befristete Angebote. Das sind Wohnungen, die f�r einen gewissen Zeitraum untervermietet werden. Hast Du 
zu Semesterbeginn noch kein Zimmer gefunden, das Dir zusagt solltest Du �berlegen, eines der befristeten Angebote wahrzunehmen. Nach 
Semesterbeginn entspannt sich n�mlich erfahrungsgem�� der Wohnungsmarkt wieder etwas. Au�erdem wird Dir die Wohungssuche vor Ort und ohne 
gro�en Zeitdruck sicher leichter fallen und Du wei�t dann auch schon etwas besser, welche Orte/Einrichtungen Du gerne in der N�he Deines 
Wohnortes h�ttest.\\
\textbf{An alle, die eine Wohnung (unter-)zuvermieten haben:}\\
Meldet euch im w�hrend unserer Gesch�ftszeiten im AStA.
\end{window}
\end{artikel}

\vspace*{-2ex}

\noindent
    \boxzt{Disclaimer}
    {
      Dieses Infoblatt soll den neuesten Stand der Dinge aus unserer Sicht verdeutlichen. Alle
      Angaben sind nach bestem Wissen und Gewissen recherchiert, trotzdem �bernehmen wir keinerlei Gew�hr.
    }
  }
  {  % rechte Spalte
\vspace*{-5ex}

\noindent
\fcolorbox{black}{light}{\begin{artikel}{Verbindungen}
Trotz der zur Zeit angespannten Wohnraumsituation solltest Du Dir gut �berlegen, ob Du ein eventuelles Angebot von Verbindungen 
wahrnehmen m�chtest. Das Angebot wird zwar vom Preis her verlockend sein, aber Du solltest Dir dar�ber im klaren sein, dass damit 
Verpflichtungen verbunden sind. N�here Infos hierzu findest Du in userem "`Denkansto� Verbindungen"'.
\end{artikel}
}
\vspace*{1ex}

\noindent
    \begin{artikel}{von Zuhause...}

...kannst Du dich auch schon nach einer Wohung umschauen, indem du dir die 
Zeitungen mit Wohnungsanzeigeteil nach Hause schicken l�sst.

Die AN (Aachener Nachrichten) und die AZ (Aachener Zeitung) haben einen kleineren Wohnungsanzeigenteil jeweils mittwochs und einen 
umfangreichen jeweils samstags. Der Anzeigenteil beider Zeitungen ist allerdings identisch. Die Bestellung der jeweiligen Ausgabe erfolgt 
schriftlich an:\\[1.5ex]
\textit{
Zeitungsverlag Aachen\\
Postfach 500110\\
52085 Aachen}\\[1.5ex]
Die AN/AZ kosten dienstags je 0,80 Euro und samstags 1 Euro. Beides zuz�glich Versandkosten von 0,35 Euro pro Ausgabe.
Au�erdem besteht die M�glichkeit sich kostenfrei im Internet den Anzeigenteil der AN anzugucken unter 
\url{http://www.aachener-nachrichten.de/}.

Das lokale Anzeigenblatt Annonce druckt jeweils dienstags und freitags Wohnungsanzeigen. Wer sich eine bestimmte Ausgabe zuschicken lassen 
m�chte, kann dies unter 0241 504021 tun. Die Annonce ist recht teuer und auch das Online-Angebot ist kostenpflichtig 
(\url{http://www.annonce.de/}).

Da zwischen Auftrag und Erhalt der Zeitung(en) wertvolle Zeit verloren geht, ist eine Suche vor Ort erfolgsversprechender. Trotzdem viel 
Erfolg beim Suchen.
 
\end{artikel}

  }
}
\end{vorderseite}
%<zweiterparameter>

\begin{rueckseite}{Wohn-Info} %Rueckseite, dass die immer so hei{\ss}t wie die Vorderseite, habe ich (noch) nicht hingekriegt.

\zweispalten{.45\linewidth}
{
  \begin{artikel}{der Mietvertrag}
Die meisten Vermieter werden Dir einen Standardmietvertrag (z.B. vom Haus- und Grundbesitzerverein) vorlegen. Legt der Vermieter einen 
eigenen (oder modifizierten) Mietvertrag vor ist auf jeden Fall Vorsicht geboten. Den Vertrag ganz genau zu lesen sollte sowieso 
selbstverst�ndlich sein. Ist dir die Bedeutung von Teilen des Vertrages unklar so solltest du fachkundigen Rat einholen (z.B. im 
AStA-Wohnungsreferat).

Unterschreibe den Vertrag auf jeden Fall erst, wenn sichergestellt ist, dass Dir nicht irgendwelche Vertragsklauseln zum Verh�ngnis werden 
k�nnten!\\

\textbf{Vorsicht vor:}
\begin{enumerate}
\item Staffelmieten (Mieterh�hungen sind bereits im Mietvertrag festgeschrieben)
\item Zeitmietvertr�gen (hier bist du einen festen Zeitraum an den Vertrag gebunden und kannst vorher quasi nicht k�ndigen)
\item�bernahme von Wohnungsinventar (Pass auf, dass der Vormieter nicht seinen Sperrm�ll an dich verscherbelt)
\end{enumerate}
\textbf{Wichtig:}
\begin{enumerate}
\item Alle Vereinbarungen und Absprachen mit dem Vermieter schriftlich festhalten. M�ndliche Absprachen sind rechtlich zwar auch bindend 
aber im Zweifelsfall deutlich schwerer einzufordern.
\item Bei der Wohungs�bergabe ein �bergabeprotokoll anfertigen. Darin alle sichtbaren M�ngel (Feuchtigkeit, Flecken im Teppich oder auf 
der 
Tapete usw.) und �bergabezustand (frisch renoviert?, wenn ja wann?) schriftlich festhalten.
\item Lass dir die Nebnkostenabrechnung des Vorjahres vorlegen, damit keine hohen Nachzahlungen f�llig werden. Der Vorauszahlungsbetrag 
kann 
unter Umst�nden sehr weit von den tats�chlichen Nebenkosten abweichen.
\item Genau notieren welche Sch�den (falls vorhanden) bis wann vom Vermieter beseitigt werden. Die Frist sollte hier genau vermerkt und 
ihre 
Einhaltung im Zweifelsfall auch eingefordert werden.
\item F�r alle F�lle solltest Du jemanden mitnehmen, der als Zeuge aussagen kann.
\end{enumerate}
\end{artikel}

}
{
  \begin{artikel}{Einwohnermeldeamt}

Hattest Du Erfolg bei der Wohnungssuche, musst Du Dich innerhalb von 14 Tagen nach Einzugsdatum (Mietvertrag) beim 
Einwohnermeldeamt 
melden, da Du sonst mit einem Bu�geld belegt werden kannst. Mitbringen musst Du dazu Deinen Personalausweis und 
den Mietvertrag. 
Es gibt zwei M�glichkeiten, Deinen Wohnsitz anzumelden: als Erst- oder als Zweitwohnsitz. 
Nach dem Meldegesetz musst Du Dich dort mit Erstwohnsitz anmelden, wo Du Deinen Lebensmittelpunkt hast. Dazu zwingen 
kann Dich jedoch Keiner, wenn Du glaubhaft darlegst, dass Du Dich quantitativ mehr am Heimatort aufh�ltst.

Die Anmeldung als Erstwohnsitz f�hrt erst einmal zu Zusatzkosten, da Du eine eigene Hausratversicherung abschlie�en 
solltest. Au�erdem kannst Du in einigen F�llen aus der Familienhaftpflichtversicherung herausfallen (in beiden F�llen im 
Versicherungsvertrag nachschauen!). Falls Du Deinen Ersatzdienst bei der Feuerwehr, beim THW o.�. ableistest, zieht ein 
Wechsel des Erstwohnsitzes au�erdem einen Wechsel Deines Arbeitsplatzes nach sich. Dein Auto musst Du ggf. in Aachen anmelden. AnwohnerInnenparkausweise werden nur an Menschen mit Erstwohnsitz in Aachen vergeben. Das gleiche gilt f�r Wohnberechtigungsscheine. Viele 
Beh�rdeng�nge k�nnen nur am Erstwohnsitz erledigt werden. Hinzu kommt, dass Du nur an Deinem Erstwohnsitz bei Kommunal- und Landtagswahlen wahlberechtigt bist. 

Von Menschen mit einem Zweitwohnsitz in Aachen wird eine Zweitwohnsitzsteuer in H�he von $10\%$ der Nettokaltmiete erhoben. 
Der Anrechnungszeitraum f�r die Steuer beginnt mit dem 
ersten Tag des Folgemonats nach Anmeldung des Zweitwohnsitzes und endet mit dem letzten Tag des Monats, in dem die 
Wohnung aufgegeben wird. Wird Dir die Wohnung aus irgendeinem Grund zu einem Preis unterhalb des Orts�blichen 
�berlassen, wird der Oberwert der Miete aus dem jeweils g�ltigen Mietspiegel zugrunde gelegt. Gemeinschaftlich genutzte 
R�ume werden anteilig auf die Fl�che der Wohnung aufgeschlagen (wichtig bei Wohngemeinschaften!). Ausnahmeregelungen f�r 
finanziell schw�chere Menschen gibt es leider nicht. 
\end{artikel}

  \begin{artikel}{Wohngeld}
Das Wohngeld stellt so eine Art
Mietzuschu� dar, der vom Sozialamt
vergeben wird. Du bekommst
es w�hrend Deines Studiums
in der Regel nur dann,
wenn Du grunds�tzlich nicht
mehr BAf�G-berechtigt bist
(z.B. wegen nicht erbrachtem
Leistungsnachweis, nicht genehmigtem
Fachrichtungswechsel
oder wegen �berschrittener
F�rderungsh�chstdauer), oder
wenn im Familienhaushalt kein
Wohnraum mehr zur Verf�gung
steht. Ein weiteres Kriterium ist
noch, da� Du Dich selbst versorgen
musst. 


N�here Informationen
dazu und zur Berechnung
der H�he des Wohngelds erh�ltst
Du in der AStA-Sozialberatung.
\end{artikel}

  \boxzt{\Large Adressen:}
{
\normalsize
\textbf{AStA}\\
Turmstr. 3, Tel. (0241) 80-93 792\\
Mo.-Fr.: 11\uuhr{}-14\Uhr\\
\textbf{Wohnheime}\\
Studentenwerk A�R,\\ Turmstr. 3,\\ Tel. (0241) 8884-401\\[1.2ex]
\parbox{.5\linewidth}{Studentenwerk der KHG e.V.,\\ Pontstr. 72,\\ Tel (0241) 4700-100}
\parbox{.5\linewidth}{Evang. StudentInnengemeinde,\\ Nizzaallee 20,\\ Tel. (0241) 918 67-0}\\[1.2ex]
\textbf{Einwohnermeldeamt Aachen Mitte}\\
Verwaltungsgeb�ude am Hauptbahnhof, Tel. (0241) 432-0\\
Mo., Di., Do.: 7\uuhr{} - 15\Uhr{}, Mi. 7\uuhr{}-15\uuuhr{}\\
\textbf{Vermittlung von WG-Zimmern}\\
HomeCompany / Mitwohnzentrale, Tel. (0241) 1 94 45
}


}
\end{rueckseite}
%<geschweifteKlammer>
\end{document}  


